% Table: drivers of the historical disaster count across six districts.
% Self-contained snippet (no preamble). Requires: booktabs, siunitx.
\begin{table}[htbp]
  \centering
  \small
  \caption{Correlation of each district attribute with the historical disaster
  count (a) and pairwise similarity of the six districts (b). Mean slope is by far
  the strongest correlate of the disaster count (Pearson $r = 0.930$, rank
  correlation $0.943$), ahead of infrastructure index ($r = -0.777$) and vegetation
  cover ($r = -0.744$), while population density is only weakly related
  ($r = 0.360$); steeper, sparsely vegetated and less developed districts therefore
  record more disasters. In (b) smaller distances mean more similar districts and
  each district's nearest neighbour is shown in bold: districts A and C form the
  most similar pair ($d = 1.59$) and E and F the next closest ($d = 2.28$), whereas
  D is the most distinctive district and E the most typical. Attributes in (a) are
  ranked by $|r|$ and the rank correlation is tied ($-0.771$) between vegetation
  cover and infrastructure index; with only six districts these values are
  indicative rather than conclusive. Source: district attribute data supplied with
  the brief.}
  \label{tab:disaster-drivers}

  \textbf{(a) Correlation with the historical disaster count (target)}\\[0.4em]
  \begin{tabular}{@{}l S[table-format=-1.3] S[table-format=-1.3] S[table-format=1.3] c@{}}
    \toprule
    Attribute (unit)                        & {$r$}  & {$\rho$} & {$|r|$} & {Rank} \\
    \midrule
    \textbf{Mean slope} ($^\circ$)          & 0.930  & 0.943    & 0.930  & \textbf{1} \\
    Infrastructure index (0--100)           & -0.777 & -0.771   & 0.777  & 2 \\
    Vegetation cover (fraction of area)     & -0.744 & -0.771   & 0.744  & 3 \\
    Population density (people\,km$^{-2}$)  & 0.360  & 0.543    & 0.360  & 4 \\
    \bottomrule
  \end{tabular}

  \par\vspace{1.1em}
  \textbf{(b) Similarity of districts (pairwise distance, smaller = more similar)}\\[0.4em]
  \begin{tabular}{@{}l r r r r r r r@{}}
    \toprule
    District (disaster count) & C & A & E & F & B & D & {Mean} \\
    \midrule
    C (8)  & --              & \textbf{1.59} & 3.44 & 4.01 & 3.50 & 6.25 & 3.76 \\
    A (12) & \textbf{1.59}   & --            & 1.91 & 2.78 & 2.65 & 4.87 & 2.76 \\
    E (16) & 3.44            & \textbf{1.91} & --   & 2.28 & 2.28 & 3.19 & 2.62 \\
    F (24) & 4.01            & 2.78          & \textbf{2.28} & -- & 3.28 & 3.75 & 3.22 \\
    B (31) & 3.50            & 2.65          & \textbf{2.28} & 3.28 & -- & 3.28 & 3.00 \\
    D (38) & 6.25            & 4.87          & \textbf{3.19} & 3.75 & 3.28 & -- & 4.27 \\
    \bottomrule
  \end{tabular}

  \par\vspace{0.5em}
  {\footnotesize
  Note: distances in (b) are Euclidean in units of pooled standard deviations over
  all five attributes, including the disaster count; the diagonal is omitted; bold
  marks each district's nearest neighbour; the right-hand column is the mean
  distance from that district to the other five. Columns in (b) follow the same
  order as the rows.}
\end{table}
